% !TeX root = ../../../main.tex
\section{Limitations}\label{sec:limitations}

The claims of this paper are bounded in ways worth listing in one place.

\begin{enumerate}
  \item \textbf{No hardware-scale validation.} The framework has been executed end-to-end only on the small reproducible instance of \cref{app:demo}, where the correct decision (\textsc{classical-only}) was known in advance; \cref{sec:audits} specifies the missing hardware-scale artifact. Until then, the framework's usefulness on genuinely uncertain instances is an argued position, not a measured one.
  \item \textbf{Screening models, not predictions.} The quantitative tools of \cref{sec:resources} are order-of-magnitude triage models whose failure modes are stated where they are introduced (see the warning box closing that section). None should be quoted as a forecast for a specific device.
  \item \textbf{Structured definitions, not theory.} The tuples $\mathcal{P}$, $\mathcal{Q}$, and $\Pi^{*}$ organize assumptions; no optimality or completeness result is derived from them.
  \item \textbf{Illustrative audits.} The genomics, scheduling, high-energy physics, and astronomy discussions audit framings; they do not report formulations, executions, or benchmarked classical baselines, and the marked TODO items in those sections are open work.
  \item \textbf{Selective related work.} The comparison in \cref{sec:related} covers the adjacent literatures the author judges most relevant, not a systematic review.
  \item \textbf{Moving baselines.} Both classical algorithms and hardware calibration data drift; every crossover estimate in this paper decays with time and must be recomputed at use.
  \item \textbf{Single-author perspective.} Prioritizations---most visibly the astronomy assessment in \cref{sec:astro}---are the author's provisional judgments under stated criteria, not community consensus.
\end{enumerate}
